% !Mode:: "TeX:UTF-8"
% !TEX program  = pdflatex
\documentclass{article}
\usepackage[ruled,vlined,linesnumbered]{algorithm2e}
\usepackage{algpseudocode}
\usepackage{amsmath}
\usepackage{amsthm}
\usepackage{graphicx}
\usepackage{subfigure}
\usepackage{float}
\usepackage{amssymb}
\usepackage{amsfonts}
\usepackage{mathrsfs}
\usepackage{pdfpages}
\usepackage{epsfig}
\usepackage{graphicx}
\usepackage{arydshln}
\usepackage{verbatim}
\usepackage{subfigure}
\usepackage{enumerate}
\usepackage{rotating}
\usepackage{threeparttable}
\usepackage{caption}
\usepackage{epsfig}
\usepackage{cite}
\usepackage{tikz}
\usetikzlibrary{shapes}
\usetikzlibrary{shapes.geometric}
\usepackage{geometry}
\geometry{a4paper, top=2.54cm, bottom=2.54cm, left=3.18cm, right=3.18cm}
\theoremstyle{definition}
\newtheorem{prob}{Problem}
\newtheorem{ans}{Answer}
\usepackage[colorlinks,linkcolor=black]{hyperref}

\linespread{1.2}

\begin{document}
	
	\section*{\LARGE [Homework 4] Poisson \& Brownian Motion}
	\section*{\large Kailing Wang 521030910356}
	
	\section*{Problem 1}
	\begin{itemize}
		\item The goal is, the first customer arrived after T-s is the last customer. The goal fails when more than 1 customers arrive in s hours. If no customer arrives, Joe will leave at T. I doubt whether the goal should be considered reached in this case, so I will consider both. The Poisson distribution with $\lambda$ is:
		\begin{equation}
			\mathbf{Pr}[X=k] = \frac{\lambda^k}{k!}e^{-\lambda} \cdot \operatorname{sign}k
		\end{equation}
		and we know we can add the $\lambda$s when considering the sum of multiple variables that satisfy Pois($\lambda$).Then $\mathbf{Pr}[X_s=1] = s\lambda e^{-s\lambda}$, $\mathbf{Pr}[X_s=0] = e^{-s\lambda}$. 
		
		Let's call the case where there must be one customer case 1, and where there can be no customer case 2. 
		
		The answer for case 1 is $s\lambda e^{-s\lambda}$ and the answer for case 2 is $(1+s\lambda)e^{-s\lambda}$.
		\item To find the optimal s, differentiate the answer in the last question with respect to s:
		
		Case 1: $\frac{d(s\lambda e^{-s\lambda})}{d s} = \lambda e^{-s\lambda}-s\lambda^2 e^{-s\lambda} = 0$, $s = \frac{1}{\lambda}$
		
		Case 2: $\frac{d((1+s\lambda) e^{-s\lambda})}{d s} = \lambda e^{-s\lambda}-s\lambda^2 e^{-s\lambda} -\lambda e^{-s\lambda} = 0$, err, I guess case 1 is correct.
	\end{itemize} 

	\section*{Problem 2}
	\begin{itemize}
		\item 
		To prove $\mathbf{Pr}[X=\lambda+k]\geq \mathbf{Pr}[X=\lambda-k-1]$, that is to prove:
		\begin{equation}
			\frac{\lambda^{\lambda+k}}{(\lambda+k)!}e^{-\lambda} \geq \frac{\lambda^{\lambda-k-1}}{(\lambda-k-1)!}e^{-\lambda}
		\end{equation}
		Simplify, 
		\begin{equation}
			\frac{\lambda^{2k+1}(\lambda-k-1)!}{(\lambda+k)!} = \frac{\lambda^{2k+1}}{\prod_{i=-k}^{k}(\lambda+i)!} \geq 1
		\end{equation}
		\begin{equation}
			\begin{aligned}
				\text{LHS} = \frac{\lambda^{2k}}{\prod_{i=1}^{k}(\lambda^2-i^2)} \geq 1
			\end{aligned}
			\label{4}
		\end{equation}
		Equation \ref{4} is obvious. Then 
		\begin{equation}
			\sum_{k = 0}^{\lambda-1}\mathbf{Pr}[X = k] \leq \sum_{k = 0}^{\lambda-1}\mathbf{Pr}[X = \lambda+k] \leq \sum_{k = 0}^{\infty}\mathbf{Pr}[X = \lambda+k] = \mathbf{Pr}[X \geq \lambda]
		\end{equation}	
		\begin{equation}
			\mathbf{Pr}[X \geq \lambda] \geq \frac{1}{2}
		\end{equation}
		\item For this problem, $\sum_{k=1}^n Y_k \sim Poisson(m\lambda)$. 
		\begin{equation}
		\begin{aligned}
			\mathbf{E}\left[f\left(Y_1, Y_2, \ldots, Y_n\right)\right] & =\sum_{k=0}^{\infty} \mathbf{E}\left[f\left(Y_1, Y_2, \ldots, Y_n\right) \mid \sum_{i=1}^n Y_i=k\right] \operatorname{Pr}\left[\sum_{i=1}^n Y_i=k\right] \\
			& \geq \sum_{k=\lambda}^{\infty}\mathbf{E}\left[f\left(Y_1, Y_2, \ldots, Y_n\right) \mid \sum_{i=1}^n Y_i=m\right] \operatorname{Pr}\left[\sum_{i=1}^n Y_i=m\right] \\
			& =\mathbf{E}\left[f\left(X_1, X_2, \ldots, X_n\right)\right]\sum_{k=\lambda}^{\infty} \operatorname{Pr}\left[\sum_{i=1}^n Y_i=m\right].
		\end{aligned}
		\end{equation}
		$\sum_{k=1}^n Y_k \sim Poisson(m\lambda)$, $\sum_{j=\lambda}^{\infty}\mathbf{Pr}[\sum_{k=1}^n Y_k = j]\geq \frac{1}{2}$
		\begin{equation}
			2 \mathbf{E}\left[f\left(Y_1, Y_2, \ldots, Y_n\right)\right] \geq \mathbf{E}\left[f\left(X_1, X_2, \ldots, X_n\right)\right].
		\end{equation}
		\item Use the corollary we get in the last question: the case there exists five students who share the same birthday can be describes as 
		\begin{equation}
			\max(X_1, \dots,  X_m) \geq 5
		\end{equation}
		where $X_i$ is the number of students whose birthday is $i$. Then we let $\sum_{i=1}^m Y_i = n$ and $Y_i\sim \operatorname{Pois}(\lambda=\frac{n}{m})$, 
		Let
		\begin{equation}
			f\left(X_1, X_2, \ldots, X_n\right) = \mathbf1 [ X \geq 5] = \mathbf1[\max(X_1, \dots,  X_m) \geq 5]
		\end{equation}
		\begin{equation}
			\begin{aligned}
				\mathbf{Pr}[\max(X_1, \dots,  X_m) \geq 5] & = \mathbf{E}\left[f\left(X_1, X_2, \ldots, X_n\right)\right] \\
				& \leq 2 \mathbf{E}\left[f\left(Y_1, Y_2, \ldots, Y_n\right)\right]\\
				& = 2\mathbf{Pr}[\max(Y_1, \dots,  Y_m) \geq 5] \\
				& = 2-2\mathbf{Pr}[\max(Y_1, \dots,  Y_m) < 5] \\
				& = 2-2(\prod_{i=1}^m\mathbf{Pr}[Y_i<5]) \\
				& \leq 2-2(0.999987)^{365} \approx 0.00947
			\end{aligned}		
		\end{equation}
		0.999987 is because $\sum_{i=0}^4\mathbf{Pr}[Y=i]\approx0.9999870714$
		
	\end{itemize}

	\section*{Problem 3}
	\begin{itemize}
		\item Prove according to definition: let's call $P(0) = c^{-\frac{1}{2}}W(ct)$
		\begin{enumerate}
			\item $P(0) = c^{-\frac{1}{2}}W(0) = 0$
			\item $\forall 0\leq t_1\leq t_2, P(t_2) - P(t_1) = c^{-\frac{1}{2}}(W(ct_2) - W(ct_1))$ is mutually independent since $W(ct_2) - W(ct_1)$ is mutually independent
			\item $\forall s,t > 0, P(s+t) - P(s) \sim \mathcal N(0,c^{-1}\cdot ct) \rightarrow \mathcal N(0,t)$
			\item $W(t)$ is continuous almost surely. So is $P(t)$
		\end{enumerate}
		Thus, $P(t)$ is a standard Brownian motion.
		\item Again, 
		\begin{enumerate}
			\item $X(0) = 0$
			\item $\forall 0\leq t_1\leq t_2, X(t_2) - X(t_1) = W(c+t_2) - W(c+t_1)$ is mutually independent
			\item $\forall s,t > 0, X(s+t) - X(s) \sim \mathcal N(0,t)$
			\item $W(t)$ is continuous almost surely. So is $X(t)$
		\end{enumerate}
		$X(t)$ is a standard Brownian motion, and
		\begin{equation}
			\begin{aligned}
				\operatorname{cov}(X(t_1), W(t_2)) & = \operatorname{cov}(W(c+t_1)-W(c), W(t_2)) \\
				& = \operatorname{cov} (W(c+t_1), W(t_2)) -\operatorname{cov} (W(c), W(t_2)) = 0
			\end{aligned}
		\end{equation}
		Thus it's independent.
		\item The original is equal to $\mathbf{Pr}[W(1)-W(\frac{1}{2})>-W(\frac{1}{2})\mid W(\frac{1}{2})>0]$, and according to the last question, $X(\frac{1}{2}) = W(\frac{1}{2}+\frac{1}{2})-W(\frac{1}{2})$ is a standard Brownian motion independent of $W(t)$. And we let $W'(t) = -W(t)$. The new problem become $\mathbf{Pr}[X(\frac{1}{2})-W'(\frac{1}{2})>0\mid W'(\frac{1}{2})<0]$.
		
		Use the total probability law, 
		\begin{equation}
			\begin{aligned}
				&\mathbf{Pr}[X(\frac{1}{2})>W'(\frac{1}{2})\mid W'(\frac{1}{2})<0] \\
				=&\mathbf{Pr}[X(\frac{1}{2})>W'(\frac{1}{2})\mid W'(\frac{1}{2})<0, X(\frac{1}{2})\geq0]\mathbf{Pr}[X(\frac{1}{2})\geq0]+\\&\mathbf{Pr}[X(\frac{1}{2})>W'(\frac{1}{2})\mid W'(\frac{1}{2})<0, X(\frac{1}{2})<0]\mathbf{Pr}[(X(\frac{1}{2}<0]\\
				=&1\cdot \frac{1}{2} + \frac{1}{2}\cdot\frac{1}{2}\\
				=&\frac{3}{4}
			\end{aligned}
		\end{equation}
		$\mathbf{Pr}[X(\frac{1}{2})>W'(\frac{1}{2})\mid W'(\frac{1}{2})<0, X(\frac{1}{2})<0] = \frac{1}{2}$ is because for two standard Brownian motion, all the situations are symmetric.
	\end{itemize}
	\section*{Problem 4}
	\begin{itemize}
		\item To prove
		\begin{equation}
			\begin{aligned}
				\mathbf{Pr}[X(t)\leq\delta] = \mathbf{Pr}\left[\xi\leq\frac{\delta-\mu\cdot t}{\sigma \sqrt{t}}\right], \\
				\mathbf{Pr}[\mu t+\sigma W(t)\leq\delta] = \mathbf{Pr}\left[\xi\leq\frac{\delta-\mu\cdot t}{\sigma \sqrt{t}}\right],
			\end{aligned}
		\end{equation}
		Here $W(t)\sim \mathcal N(0, t)$. View $W(t)$ as $\sqrt{t}\xi$, 
		\begin{equation}
			\begin{aligned}
				\mathbf{Pr}[\mu t+\sigma\sqrt{t} \xi\leq\delta] = \mathbf{Pr}\left[\xi\leq\frac{\delta-\mu\cdot t}{\sigma \sqrt{t}}\right],
			\end{aligned}
		\end{equation}
		This is much too obvious.
		\item According to the last question, 
		\begin{equation}
			\begin{aligned}
				\mathbf{Pr}[X(t)\geq0] = \mathbf{Pr}\left[\xi\geq\frac{-\mu\cdot t}{\sigma \sqrt{t}}\right],\\
				\mathbf{Pr}[X(t)\leq\delta] = \mathbf{Pr}\left[\xi\leq\frac{\delta-\mu\cdot t}{\sigma \sqrt{t}}\right], \\
			\end{aligned}
		\end{equation}
		Let a new $\xi$ be $-\xi$ and replace it in the equations above:
		\begin{equation}
			\begin{aligned}
				\mathbf{Pr}[X(t)\geq0] &= \mathbf{Pr}\left[\xi\geq\frac{-\mu\cdot t}{\sigma \sqrt{t}}\right],\\
				\mathbf{Pr}[0\leq X(t)\leq\delta] &= \mathbf{Pr}\left[\frac{\mu\cdot t-\delta}{\sigma \sqrt{t}}\leq\xi\leq\frac{\mu\cdot t}{\sigma \sqrt{t}} \right], \\
				\int_0^\infty\mathbf{Pr}[0\leq X(t)\leq\delta]dt &= \int_0^\infty\mathbf{Pr}\left[\frac{\mu\cdot t-\delta}{\sigma \sqrt{t}}\leq\xi\leq\frac{\mu\cdot t}{\sigma \sqrt{t}} \right]dt \\
			\end{aligned}
		\end{equation}
		And the LHS is the definition of the expectation of T.
		\item This is to find $\mathbf{Pr}[0\leq X(t)\leq\delta] = \mathbf{Pr}[f(0,\xi)\leq t \leq f(\delta, \xi)]$. $X(t) = \mu t+\sigma\sqrt{t} \xi$. Simply we need to solve the quadratic equation.
		\begin{equation}
			\begin{aligned}
				t&\geq \frac{\sigma^2\xi^2}{\mu^2} \\
				\sqrt{t} &\leq \frac{-\sigma\xi\pm\sqrt{\sigma^2\xi^2+4\mu\delta}}{2\mu}\\
				f(0, \xi) &= \frac{\sigma^2\xi^2}{\mu^2}\\
				f(\delta, \xi) &= \frac{\delta}{\mu} + \frac{\sigma^2\xi^2}{2\mu^2} + \frac{\sigma\xi\sqrt{\sigma^2\xi^2+4\mu\delta}}{2\mu^2}
			\end{aligned}
		\end{equation}
		\item $\mathbf E{f(\delta,\xi)} = \mathbf E[\frac{\delta}{\mu} + \frac{\sigma^2\xi^2}{2\mu^2} + \frac{\sigma\xi\sqrt{\sigma^2\xi^2+4\mu\delta}}{2\mu^2}]=\frac{\delta}{\mu} + \frac{\sigma^2\mathbf E[\xi^2]}{2\mu^2} + \frac{\sigma\mathbf E[\xi\sqrt{\sigma^2\xi^2+4\mu\delta}]}{2\mu^2}$
		
		$\mathbf E[\xi^2] = 1$, $\mathbf E[\xi\sqrt{\sigma^2\xi^2+4\mu\delta}] = 0$
		
		$\mathbf E[f(\delta,\xi)] = \frac{\delta}{\mu}+\frac{\sigma^2}{2\mu^2}$ 
		\item $\mathbf E [T] = \int_0^\infty\mathbf{Pr}\left[\frac{\mu\cdot t-\delta}{\sigma \sqrt{t}}\leq\xi\leq\frac{\mu\cdot t}{\sigma \sqrt{t}} \right]dt$
		
		$=\int_0^\infty\mathbf{Pr}[f(0,\xi)\leq t \leq f(\delta, \xi)]dt$
		
		$=\int_0^\infty\mathbf{Pr}[t \leq f(\delta, \xi)]dt-\int_0^\infty\mathbf{Pr}[t \leq f(0, \xi)]dt$
		
		$=\frac{\delta}{\mu}+\frac{\sigma^2}{2\mu^2}-\frac{\sigma^2}{2\mu^2} = \frac{\delta}{\mu}$
	\end{itemize}
	\begin{thebibliography}{99}
		\bibitem{zhangsp2023} C. Zhang, Stochastic Processes (Spring 2023) [Online]. Available: \url{http://chihaozhang.com/teaching/SP2023/} [Accessed Mar. 24, 2023].
	\end{thebibliography}
\end{document}

